\begin{abstract}
Malicious code poisoning in the AI model supply chain has emerged as a critical threat. Attackers embed hidden payloads into model artifacts distributed through open-source hubs, enabling arbitrary code execution upon deserialization or inference. Existing defenses rely on static signature scanning and prior knowledge of attack patterns, making them inherently vulnerable to evasion and zero-day exploits.
We observe that benign AI models exhibit highly regular runtime execution templates, while malicious models inevitably deviate from these patterns. Building on this observation, we propose \toolname{}, a dynamic defense that requires no prior knowledge of attack methodologies. \toolname{} combines kernel-level system call monitoring with Python-level semantic context and bridges the semantic gap between low-level runtime evidence and high-level model intent through cross-layer analysis.
We evaluate \toolname{} on existing MalHug dataset and Advanced TensorAbuse Synthetic dataset spanning four attack categories. \toolname{} achieves an F1-score of 0.9870 on the MalHug and 1.0000 on the Advanced TensorAbuse Synthetic Dataset, outperforming all baseline methods. \toolname{} introduces only 5.69\% average runtime overhead with 9.79 seconds detection latency, representing an acceptable cost for pre-deployment model vetting. 
\end{abstract}
